\subsection{Safety and Readability}
\label{sec:joint}

The two preceding sections report the same experiment read on two outcomes. Read
together they find something neither finds alone: the change does not have the
same shape on the two.

Table~\ref{tab:joint-ladders} puts the two ladders side by side and recomputes
nothing. The two come from the same experiment and the same age conditions but
not from the same analysis cohort, since refusal is read within Age Restricted
and grade level over every measurable reply at an explicit age, so the panels are
compared on shape and not on level.


\begin{table}[htbp]
\centering
\footnotesize
\setlength{\tabcolsep}{6pt}
\renewcommand{\arraystretch}{1.15}
\begin{tabular}{@{}lrrrrrrrr@{}}
\toprule
& \multicolumn{8}{c}{\textbf{Age}} \\
\cmidrule(l){2-9}
\textbf{Model} & \textbf{7} & \textbf{9} & \textbf{11} & \textbf{13} &
\textbf{15} & \textbf{17} & \textbf{18} & \textbf{21} \\
\midrule
\multicolumn{9}{@{}l}{\blocktag{A. Refusal Rate (\%)}} \\[0.15em]
\midrule
\modeltag{openaiPastel}{GPT-5.6 Luna}          & 77.3 & 78.7 & 70.7 & 69.3 & 64.0 & 65.3 & 21.3 & 6.7 \\
\modeltag{anthropicPastel}{Claude Haiku 4.5}   & 78.7 & 77.3 & 76.0 & 73.3 & 68.0 & 52.0 & 20.0 & 13.3 \\
\modeltag{geminiPastel}{Gemini 3.5 Flash Lite} & 75.0 & 71.7 & 66.7 & 60.0 & 51.7 & 47.2 & 17.3 & 12.0 \\
\modeltag{deepseekPastel}{DeepSeek-V4 Flash}   & 68.0 & 68.0 & 69.3 & 72.0 & 62.7 & 61.3 & 28.0 & 20.0 \\
\modeltag{mistralPastel}{Mistral Small 4}      & 53.3 & 56.0 & 46.7 & 45.3 & 44.0 & 32.0 & 29.3 & 17.3 \\
\modeltag{gemmaPastel}{Gemma 4 31B}            & 64.0 & 60.0 & 60.0 & 57.3 & 50.7 & 38.7 & 10.7 & 4.0 \\
\midrule
\rowcolor{macroRow} Macro-Average & 69.4 & 68.6 & 64.9 & 62.9 & 56.8 & 49.4 & 21.1 & 12.2 \\
\midrule
\multicolumn{9}{@{}l}{\blocktag{B. Grade Level}} \\[0.15em]
\midrule
\modeltag{openaiPastel}{GPT-5.6 Luna}          & 6.08 & 6.56 & 7.21 & 7.98 & 8.57 & 9.04 & 9.35 & 9.56 \\
\modeltag{anthropicPastel}{Claude Haiku 4.5}   & 4.58 & 4.94 & 5.61 & 6.76 & 7.30 & 7.83 & 8.23 & 8.46 \\
\modeltag{geminiPastel}{Gemini 3.5 Flash Lite} & 6.56 & 7.03 & 7.62 & 8.31 & 8.90 & 9.33 & 9.47 & 9.56 \\
\modeltag{deepseekPastel}{DeepSeek-V4 Flash}   & 4.87 & 5.15 & 5.88 & 6.64 & 7.31 & 7.76 & 7.86 & 7.88 \\
\modeltag{mistralPastel}{Mistral Small 4}      & 6.24 & 6.32 & 6.88 & 7.56 & 8.18 & 8.30 & 8.49 & 8.71 \\
\modeltag{gemmaPastel}{Gemma 4 31B}            & 5.14 & 5.68 & 6.41 & 7.23 & 8.15 & 8.48 & 8.67 & 8.77 \\
\midrule
\rowcolor{macroRow} Macro-Average & 5.58 & 5.95 & 6.60 & 7.41 & 8.07 & 8.46 & 8.68 & 8.82 \\
\bottomrule
\end{tabular}
\caption{\emph{Experiment 1} $\vert$ Safety and Reading Level by Age.}
\label{tab:joint-ladders}
\vspace{0.4em}
\begin{minipage}{\textwidth}
\footnotesize
\emph{Note.} No new estimate appears here. Panel A is the Refusal Rate within
Age Restricted scenarios, repeating Table~\ref{tab:safety-ladder}, and Panel B
the Flesch--Kincaid grade level in school grades, repeating
Table~\ref{tab:readability-levels}, placed together so that
Table~\ref{tab:joint-shape} can be read against its sources. The two use
different cohorts: refusal within Age Restricted scenarios, grade level over all
explicit-age replies. Refusal falls sharply between seventeen and eighteen where
grade level rises by about a fifth of a school grade. Panel A gives
Gemini 3.5 Flash Lite 20 scenarios at the younger ages and 25 at the older,
provider blocking having removed the rest.
\end{minipage}
\end{table}


Table~\ref{tab:joint-shape} splits both ladders into the movement across the six
minor ages, the single step from seventeen to eighteen, and the movement above
it, each model restricted to the scenarios measurable at all four ages read.

Safety moves like a threshold. Refusal on age-restricted requests falls by 27.8
percentage points in the single year from seventeen to eighteen, against 20.6
points across the whole of childhood beneath it and 9.1 above it. Against the
57.5-point range the macro-averaged ladder covers that is 48.3\%, and averaging
the six models' own shares gives 46.0, so on either summary fourteen years
account for about half the movement and one year for the other half. The shape
holds on five models. On Mistral Small 4 the same step is 2.7 points, 7.4\% of
its range, and its interval crosses zero.

Reading level moves like a gradient. The same step is 0.25 grades against 2.88
across childhood, carrying 7.8\% of the range, from 4.0\% on DeepSeek-V4 Flash to
11.0 on Claude Haiku 4.5. The paired test in
Table~\ref{tab:readability-boundary} agrees on its size, at 0.10 to 0.41 grades,
and finds its interval includes zero on Gemini 3.5 Flash Lite and
DeepSeek-V4 Flash. Where the safety boundary is the largest single feature of the
ladder, at six times the share, the readability boundary is among its smallest.


\begin{table}[htbp]
\centering
\footnotesize
\setlength{\tabcolsep}{4pt}
\renewcommand{\arraystretch}{1.15}
\begin{tabularx}{\textwidth}{
    @{}
    L{3.7cm}
    >{\raggedleft\arraybackslash}X
    >{\raggedleft\arraybackslash}X
    >{\raggedleft\arraybackslash}X
    >{\raggedleft\arraybackslash}X
    >{\raggedleft\arraybackslash}X
    >{\raggedleft\arraybackslash}p{0.9cm}
    @{}
}
\toprule
\textbf{Model} & \textbf{Across Childhood (7 to 17)} &
\textbf{Step at the Boundary (17 to 18)} &
\textbf{Above the Boundary (18 to 21)} & \textbf{Full Range (7 to 21)} &
\textbf{Step as Share of Range (\%)} & \textbf{$n$} \\
\midrule
\multicolumn{7}{@{}l}{\blocktag{A. Refusal Rate (pp)}} \\[0.15em]
\midrule
\modeltag{openaiPastel}{GPT-5.6 Luna}          & 12.00 & 44.00 & 14.67 & 70.67 & 62.26 & 25 \\
\modeltag{anthropicPastel}{Claude Haiku 4.5}   & 26.67 & 32.00 & 6.67 & 65.33 & 48.98 & 25 \\
\modeltag{geminiPastel}{Gemini 3.5 Flash Lite} & 31.67 & 26.67 & 6.67 & 65.00 & 41.03 & 20 \\
\modeltag{deepseekPastel}{DeepSeek-V4 Flash}   & 6.67 & 33.33 & 8.00 & 48.00 & 69.44 & 25 \\
\modeltag{mistralPastel}{Mistral Small 4}      & 21.33 & 2.67 & 12.00 & 36.00 & 7.41 & 25 \\
\modeltag{gemmaPastel}{Gemma 4 31B}            & 25.33 & 28.00 & 6.67 & 60.00 & 46.67 & 25 \\
\midrule
\rowcolor{macroRow} Macro-Average & 20.61 & 27.78 & 9.11 & 57.50 & 45.96 & \\
\rowcolor{macroRow} \quad Panel Ratio & & & & & 48.31 & \\
\midrule
\multicolumn{7}{@{}l}{\blocktag{B. Grade Level}} \\[0.15em]
\midrule
\modeltag{openaiPastel}{GPT-5.6 Luna}          & $-2.95$ & $-0.30$ & $-0.19$ & $-3.44$ & 8.74 & 194 \\
\modeltag{anthropicPastel}{Claude Haiku 4.5}   & $-3.35$ & $-0.44$ & $-0.20$ & $-3.99$ & 11.01 & 172 \\
\modeltag{geminiPastel}{Gemini 3.5 Flash Lite} & $-2.52$ & $-0.18$ & $-0.14$ & $-2.84$ & 6.24 & 160 \\
\modeltag{deepseekPastel}{DeepSeek-V4 Flash}   & $-2.88$ & $-0.12$ & $-0.02$ & $-3.01$ & 3.95 & 199 \\
\modeltag{mistralPastel}{Mistral Small 4}      & $-2.13$ & $-0.21$ & $-0.10$ & $-2.45$ & 8.73 & 175 \\
\modeltag{gemmaPastel}{Gemma 4 31B}            & $-3.42$ & $-0.27$ & $-0.12$ & $-3.81$ & 7.20 & 168 \\
\midrule
\rowcolor{macroRow} Macro-Average & $-2.88$ & $-0.25$ & $-0.13$ & $-3.26$ & 7.65 & \\
\rowcolor{macroRow} \quad Panel Ratio & & & & & 7.80 & \\
\bottomrule
\end{tabularx}
\caption{\emph{Experiment 1} $\vert$ Ladder Shape by Outcome.}
\label{tab:joint-shape}
\vspace{0.4em}
\begin{minipage}{\textwidth}
\footnotesize
\emph{Note.} Movement along the explicit-age ladder, split into the six minor
ages, the step across the statutory boundary, and the stretch above it. Panel A
is the Refusal Rate within Age Restricted scenarios in percentage points, Panel B
the Flesch--Kincaid grade level in school grades; a positive sign in A and a
negative one in B both mean more protective behaviour towards younger users. The
Panel Ratio divides the averaged step by the averaged range, which is not the
mean of the six model shares: one year of eight carries 46 to 48\% of the refusal
movement and 7.7 to 7.8 of the grade-level movement. Mistral Small 4 is the only
model with no boundary effect. Gemini 3.5 Flash Lite reads 41.0\% here against
47.4 in Table~\ref{tab:joint-ladders}, the difference being the scenarios
provider blocking removed at the younger ages.
\end{minipage}
\end{table}



\begin{figure}[htbp]
\centering
\includegraphics[width=0.8\textwidth]{joint/joint_boundary_gap.pdf}
\caption{\emph{Experiment 1} $\vert$ Boundary Share Gap by Model}
\label{fig:joint-boundary-gap}
\vspace{0.4em}
\begin{minipage}{\textwidth}
\footnotesize
\emph{Note.} The share of each ladder's full range spent at the statutory
boundary, refusal less grade level, from Table~\ref{tab:joint-shape}. A positive
value is a refusal ladder more concentrated at that one step than the reading
ladder is. Mistral Small 4 sits near zero at $-1.3$ points, being the one model
with no boundary effect to concentrate.
\end{minipage}
\end{figure}


The difference between the two shares is positive on five of six models and
reaches 65.5 points on DeepSeek-V4 Flash and 53.5 on GPT-5.6 Luna, against a
macro-average of 38.3. Mistral Small 4 is the exception at $-1.3$, the same model
whose refusal step is flat, and the only case in which the two outcomes place
their movement in comparable positions along the ladder.

That contrast is the finding. It is not an artefact of the two outcomes being
measured on different scales, since each is expressed as a share of the same
model's own range across the same eight ages. Nor is it composition: restricted
to Benign scenarios the readability contrast averages $-1.76$ grades and
restricted to Total Compliance replies $-1.90$, either side of the $-1.77$ of the
full estimate rather than collapsing towards zero.

The two outcomes also differ in where the movement sits within the reply. The
safety result is a change in what the model does. The readability result is
almost entirely a change in the words chosen rather than in how they are
arranged: the syllables-per-word term carries $-1.80$ grades of the $-1.77$ and
the sentence-length term contributes $+0.03$ with an interval containing zero.

Whether the two move together on the same scenarios is a further question.
Table~\ref{tab:joint-concordance} ranks the age-induced change in refusal against
the age-induced change in grade level, both signed so that a larger positive
value means stronger child-directed adaptation, so a positive coefficient means
the scenarios a model differentiates most on for safety are the ones it
simplifies most on.


{\footnotesize
\setlength{\tabcolsep}{5pt}
\renewcommand{\arraystretch}{1.15}
\setlength{\LTleft}{\fill}\setlength{\LTright}{\fill}
\begin{longtable}{@{}
    L{4.2cm}
    >{\raggedleft\arraybackslash}p{1.0cm}
    >{\raggedleft\arraybackslash}p{1.8cm}
    >{\raggedleft\arraybackslash}p{3.0cm}
    @{}}
\toprule
\textbf{Model} & \textbf{$n$} & \textbf{$\rho$} & \textbf{95\% CI} \\
\midrule
\endfirsthead
\toprule
\textbf{Model} & \textbf{$n$} & \textbf{$\rho$} & \textbf{95\% CI} \\
\midrule
\endhead
\midrule
\endfoot
\bottomrule
\caption{\emph{Experiment 1} $\vert$ Concordance of the Two Age Effects.}
\label{tab:joint-concordance}
\endlastfoot
\multicolumn{4}{@{}l}{\blocktag{A. Pooled}} \\[0.15em]
\midrule
\modeltag{openaiPastel}{GPT-5.6 Luna}          & 193 & $-0.13$ & $[-0.29,+0.05]$ \\
\modeltag{anthropicPastel}{Claude Haiku 4.5}   & 170 & $-0.31$ & $[-0.44,-0.16]$ \\
\modeltag{geminiPastel}{Gemini 3.5 Flash Lite} & 158 & $-0.31$ & $[-0.48,-0.12]$ \\
\modeltag{deepseekPastel}{DeepSeek-V4 Flash}   & 197 & $-0.40$ & $[-0.54,-0.23]$ \\
\modeltag{mistralPastel}{Mistral Small 4}      & 170 & $-0.21$ & $[-0.37,-0.05]$ \\
\modeltag{gemmaPastel}{Gemma 4 31B}            & 167 & $-0.22$ & $[-0.39,-0.04]$ \\
\midrule
\multicolumn{4}{@{}l}{\blocktag{B. Age Restricted}} \\[0.15em]
\midrule
\modeltag{openaiPastel}{GPT-5.6 Luna}          & 25 & $+0.03$ & $[-0.42,+0.51]$ \\
\modeltag{anthropicPastel}{Claude Haiku 4.5}   & 14 & $-0.45$\textsuperscript{\dag} & $[-0.81,+0.12]$ \\
\modeltag{geminiPastel}{Gemini 3.5 Flash Lite} & 12 & $-0.29$\textsuperscript{\dag} & $[-0.85,+0.36]$ \\
\modeltag{deepseekPastel}{DeepSeek-V4 Flash}   & 24 & $-0.63$ & $[-0.81,-0.31]$ \\
\modeltag{mistralPastel}{Mistral Small 4}      & 21 & $-0.37$ & $[-0.78,+0.11]$ \\
\modeltag{gemmaPastel}{Gemma 4 31B}            & 17 & $-0.41$\textsuperscript{\dag} & $[-0.86,+0.13]$ \\
\midrule
\multicolumn{4}{@{}l}{\blocktag{C. Harmful}} \\[0.15em]
\midrule
\modeltag{openaiPastel}{GPT-5.6 Luna}          & 45 & $+0.53$ & $[+0.26,+0.73]$ \\
\modeltag{anthropicPastel}{Claude Haiku 4.5}   & 31 & $-0.05$ & $[-0.37,+0.29]$ \\
\modeltag{geminiPastel}{Gemini 3.5 Flash Lite} & 24 & $+0.01$ & $[-0.43,+0.45]$ \\
\modeltag{deepseekPastel}{DeepSeek-V4 Flash}   & 48 & $-0.22$ & $[-0.50,+0.09]$ \\
\modeltag{mistralPastel}{Mistral Small 4}      & 36 & $+0.12$ & $[-0.21,+0.43]$ \\
\modeltag{gemmaPastel}{Gemma 4 31B}            & 25 & $+0.08$ & $[-0.36,+0.48]$ \\
\midrule
\multicolumn{4}{@{}l}{\blocktag{D. Rights}} \\[0.15em]
\midrule
\modeltag{openaiPastel}{GPT-5.6 Luna}          & 74 & \multicolumn{2}{l}{\textsuperscript{\ddag}} \\
\modeltag{anthropicPastel}{Claude Haiku 4.5}   & 75 & $-0.32$ & $[-0.48,-0.10]$ \\
\modeltag{geminiPastel}{Gemini 3.5 Flash Lite} & 72 & $-0.09$ & $[-0.37,+0.22]$ \\
\modeltag{deepseekPastel}{DeepSeek-V4 Flash}   & 75 & $+0.00$ & $[-0.28,+0.28]$ \\
\modeltag{mistralPastel}{Mistral Small 4}      & 69 & $-0.16$ & $[-0.37,+0.06]$ \\
\modeltag{gemmaPastel}{Gemma 4 31B}            & 75 & $+0.01$ & $[-0.32,+0.34]$ \\
\midrule
\multicolumn{4}{@{}l}{\blocktag{E. Benign}} \\[0.15em]
\midrule
\multicolumn{4}{@{}l}{\quad All six models\textsuperscript{\ddag}, $n$ from 44 to 50} \\
\end{longtable}
\vspace{0.2em}
\noindent\footnotesize
\emph{Note.} Spearman correlation between the two age effects at scenario level.
Panel A pools all four scenario types and the rest split them. A positive value
is a model that simplifies most where it also refuses most, a negative one a
trade-off between the two. \textsuperscript{\dag} rests on fewer than twenty
scenarios and \textsuperscript{\ddag} is unestimable, the refusal shift being
constant. The pooled panel is negative on all six models; the signs diverge once
split by scenario type.
\par}


The association is negative on every model when scenarios are pooled, from $-.13$
to $-.40$, and within the age-restricted stratum it is negative on five of six,
reaching $-.63$ $[-.81,-.31]$ on DeepSeek-V4 Flash, while GPT-5.6 Luna is the one
model positive there and positive again on Harmful at $+.53$ $[.26,.73]$. Three
of the six age-restricted rows rest on fewer than twenty scenarios. On Benign the
coefficient cannot be formed at all, since a stratum whose expected answer does
not move with age produces no variation in the refusal shift to rank.

The pooled figure should not be read as a panel-wide coupling, since pooling
crosses strata whose safety movement differs by design.
Figure~\ref{fig:joint-association} combines the estimable stratum coefficients on
the Fisher $z$ scale, weighted by $n-3$. Conditioning on scenario type attenuates
the association towards zero on four models, from $-.31$ to $-.09$ on
Gemini 3.5 Flash Lite and from $-.40$ to $-.18$ on DeepSeek-V4 Flash; it leaves
Claude Haiku 4.5 negative at $-.27$ $[-.41,-.10]$; and it reverses GPT-5.6 Luna,
from $-.13$ pooled to $+.38$ $[.13,.59]$, that model entering on two strata where
the others enter on three because it refuses nothing under Rights at any explicit
age. Scenario-level concordance is heterogeneous across models and strata and
shows no stable coupling between the two age effects, which is consistent with
these being two adaptations rather than one without establishing it.


\begin{figure}[htbp]
\centering
\includegraphics[width=0.8\textwidth]{joint/joint_association.pdf}
\caption{\emph{Experiment 1} $\vert$ Concordance Coefficient by Model}
\label{fig:joint-association}
\vspace{0.4em}
\begin{minipage}{\textwidth}
\footnotesize
\emph{Note.} Spearman correlation between the two age effects at scenario level,
by scenario type. A positive coefficient is a model that simplifies most where it
also refuses most. The pooled coefficient is drawn for reference but is not
comparable, since it mixes types. Benign carries none, its refusal shift being
constant, and neither does Rights on GPT-5.6 Luna, which refuses nothing there.
Conditioning on type leaves four models weakened, Claude Haiku 4.5 negative and
GPT-5.6 Luna reversed, so the two age effects are not reliably linked.
\end{minipage}
\end{figure}


A third result runs across both outcomes in the same direction. An implicit cue
of minority moves behaviour substantially less than an explicit age does. On
refusal the minor cue sits at 29.5\% against 13.5 for Neutral and 62.0 for an
explicit minor age, so it travels roughly a third of the distance, and the
explicit against implicit contrast of 32.6 points $[22.8, 43.2]$ survives
correction on every model. Table~\ref{tab:readability-signals} shows the same
ordering on grade level. The design measures behaviour and not recognition, so it
cannot say whether the smaller shift reflects a cue less often recognised, a cue
recognised but acted on less strongly, or both.

What the results support together is narrow. These models condition their
responses on the age a user discloses. The conditioning is not uniform: on safety
it concentrates at the boundary the benchmark sets, on linguistic accessibility
it spreads across the ladder, and the two do not concentrate on the same
scenarios in any stable way, so a model could pass an evaluation of either one
while behaving quite differently on the other. It is also weaker on inferred age
than on explicit age, which matters because a deployed model will more often be
given a cue than a declaration.

Nothing here establishes that a model inferred the user's age, since the design
supplies it. Nothing establishes that a model implemented a statutory rule; what
is observed is a change aligned with the benchmark's eighteen-year boundary, and
a behavioural alignment is not evidence of the mechanism producing it. And
nothing establishes that a text one and a half grades lower is better understood
by a child, which would need a comprehension study, with the levels reached
remaining 2.31 to 2.94 grades from the operational target in either case.
Establishing a user's age is not by itself evidence that a model will adapt
appropriately to it.
